\reviewsection{Dignity Changes the Outcome We Count}

Dunlap, Strain, and Fox place positive behavior support inside quality of life, social validity, ecological validity, collaboration, prevention, and sustainable systems change \citep{dunlap2012pbs}. Those commitments matter because technical effectiveness is not the same as dignity. A quieter student may make a classroom easier for an adult without becoming safer, more understood, or more connected. A socially meaningful outcome asks whether the student can communicate, enter learning, protect a boundary, build relationships, recover after difficulty, and exercise greater control over everyday life. Next, Gardner's first-hand report makes the hidden outcomes of intervention impossible to dismiss \citep{gardner2017firsthand}. Autistic people and people with other developmental disabilities described harm when intervention centered on appearing less disabled, suppressing harmless regulation, or teaching unquestioned compliance. They valued self-understanding, sensory regulation, autonomy, cultural competence, trauma-sensitive care, inclusion, and self-advocacy. Their accounts show why lived experience isn't an emotional supplement added after professional data. It can reveal shame, masking, weakened boundaries, trust loss, or a hidden, more harmful replacement behavior that a frequency count misses.

Dignity therefore changes what I record. ``The student refused'' is shallow evidence. Richer evidence includes what happened before the refusal, whether the task and transition were understandable, what sensory and social demands were present, which communication routes were available, how adults responded, what the student indicated, what family knowledge adds, and what changed when support changed. Student behavior belongs in the record, but so do adult behavior and environmental design.

\reviewsection{Consent Must Survive the Moment of Difficulty}

\vspace{-0.2cm}
Consent cannot exist only when a student agrees with the adult. Gardner's participants warn that repeatedly overriding ``no'' can weaken a person's ability to set boundaries and protect themselves \citep{gardner2017firsthand}. In school, consent does not mean adults abandon instruction, safety, or accountability. It means that ``no,'' ``stop,'' ``not yet,'' silence, movement away, and other forms of hesitation are treated as information before they are treated as defiance. Jorgensen makes that obligation concrete by arguing that presuming competence is only the beginning; adults must construct communicative competence through age-respectful instruction, AAC and other communication supports, collaboration, and access throughout the day \citep{jorgensen2018competence}. A system limited to requesting preferred items does not give a learner enough language for disagreement, academic thinking, social connection, refusal, repair, or self-advocacy. Speech, typing, drawing, pointing, gesture, AAC, a familiar device, partner support, silence, and delayed response can all carry meaningful evidence. Consent can therefore be designed into ordinary routines. Students can know the purpose and sequence of a support, choose among equivalent participation routes, ask for another tool or more time, pass and re-enter through a visible next step, and decide whether reflection stays private or is shared. When immediate danger exists, adults must protect people using the least intrusive effective response. That responsibility does not turn every movement, refusal, stim, delay, or unconventional communication into a safety crisis.

\vspace{-0.2cm}
\reviewsection{Neurodiversity Changes the Purpose of Support}

\vspace{-0.2cm}
Respecting neurodiversity does not mean pretending that every condition is harmless or that teaching is unnecessary. It means refusing to define success as looking neurotypical. Movement may support attention. A familiar object may support predictability. An indirect or delayed response may contain careful thinking. A learner may need a safer or more effective communication pathway, but that pathway does not need to imitate the adult's preferred style. This is where Module 1 and Module 2 meet. Diagnostic knowledge can open services, and explicit instruction can expand access, but neither should become permission to rank bodyminds by normality. The question is whether it gives the student a more dependable way to learn, communicate, participate, maintain boundaries, and belong without surrendering identity.
